\subsection{トピックと参考資料の対応}

この表は、SWEBOK の「トピック対参考資料の行列」を学習用に簡略化したものである。詳しく学ぶときは、各トピックに対応する文献を読む。

\par\smallskip\noindent\refstepcounter{table}\label{tab:ch15-04}表 \thetable: トピックと参考資料の対応の整理\par\smallskip
表 \thetable は、トピックと参考資料の対応の整理を比較・確認しやすい形で整理したものである。\par\smallskip
\begin{longtable}{L{0.42\linewidth} L{0.53\linewidth}}
\toprule
トピック & 主な参考資料 \\
\midrule
ソフトウェア工学経済の基礎 & Tockey, \textit{Return on Software}; Engineering Economy \\
提案、キャッシュフロー、時間価値、等価性、比較基準、代替案 & Tockey, \textit{Return on Software} \\
無形資産とビジネスモデル & Sanchez-Segura らの無形資産研究、Business Model 関連文献 \\
工学的意思決定プロセス & Tockey, \textit{Return on Software}; ISO/IEC/IEEE 12207; ISO/IEC/IEEE 15288 \\
営利組織の意思決定 & Tockey, \textit{Return on Software} \\
置換判断、撤退判断 & Tockey; ISO/IEC/IEEE 15288 \\
非営利組織の意思決定 & Tockey, \textit{Return on Software} \\
現在時点の経済的意思決定 & Tockey, \textit{Return on Software} \\
複数属性意思決定 & Tockey; Gilb; ATAM \\
無形資産の識別と特徴付け & SIPAC 関連文献、Value Proposition Design \\
見積り & Tockey; McConnell, \textit{Software Estimation}; Stutzke \\
実務上の考慮 & PMBOK Guide; Software Extension to PMBOK; systems thinking 文献 \\
関連概念 & Sommerville; Fairley; PMBOK Guide \\
\bottomrule
\end{longtable}
